\section{Parameters and Integration}
\label{sec:parameters}

This section explains how parameters interact across the ARC-SPRUCE stack: the lattice signature layer, the rational
hash, the SmallWood proof system, and the PRF. Detailed parameter tables and extended benchmark data are deferred to
Appendix~\ref{app:ext-params}.

\subsection{Parameter dependencies}
All components share the same base ring \(\Rq\) and modulus \(q\):
\begin{itemize}[leftmargin=1.25em]
  \item The signature map \(A\in\Rq^{n\times m}\) and hash key \(B\) are defined over \(\Rq\).
  \item SmallWood constraints are evaluated over the base field \(\Fq\) and therefore use the same modulus \(q\).
  \item The PRF is instantiated over \(\Fq\) so that its constraints are native to the same field.
\end{itemize}
The coefficient bound \(\BoundMsg\) governs the sampling ranges for the message and bounded randomness and the
membership polynomials used by the NIZKs. For the aligned commitment and hiding theorem, the linear-side distributions
\(\chi_s,\chi_e\) are instantiated as bounded-support or truncated distributions inside
\([-\BoundMsg,\BoundMsg]\).

\subsection{Sampler parameters}
The trapdoor sampler must output short signatures \(u\) with overwhelming probability and with a distribution close to
a discrete Gaussian over the appropriate lattice coset. This requires choosing:
\begin{itemize}[leftmargin=1.25em]
  \item ring degree \(N\) and modulus \(q\),
  \item trapdoor quality,
  \item Gaussian widths for the sampler, and
  \item an acceptance predicate that yields a public verification bound \(\BoundSig\).
\end{itemize}
We use an NTRU-style trapdoor together with a Gaussian preimage sampler; details and tuning rules are given in
Appendix~\ref{app:gaussian-sampler}.

\subsection{vSIS and rational-hash parameters}
\paragraph{Rational-hash admissibility.}
The BB-tran target is evaluated only when \(B_3-R_1\in \Rq^\times\). In the issuance model, this admissibility
condition is enforced during honest target derivation by restart. The compiled pre-sign proof of
Section~\ref{subsec:presign-compiled} certifies the aligned public identity
\[
T = B_0 + U\vecc + Z,
\qquad
(B_3-R_1)\Had Z = 1,
\]
while the compiled post-sign proof certifies the eliminated coefficient-domain relation
\[
(B_3-R_1)\Had\bigl(T-L_B(m,R_0)\bigr)=1.
\]

\paragraph{No wrap-around in bounds.}
Membership checks for \([ -\BoundMsg,\BoundMsg]\) are implemented via vanishing polynomials over \(\Fq\), so one
requires \(\BoundMsg \ll q\) to avoid modular wrap-around artifacts.

\paragraph{Non-blind reduction parameters.}
The non-blind reduction imported from DKLW25 is parameterized by a norm bound
\(\beta_{\mathrm{DKLW}}\) in the norm used in \cite{EPRINT:DKLW25}, by the inadmissibility bound \(\delta\), and by
the predicate condition \(P\). For the present instantiation, the external reduction basis is the BB-tran /
\(\mathsf{GenISIS}_f\) / \(\mathsf{IntGenISIS}_f\) line summarized in Lemma~\ref{lem:vsis-unforgeability}, which
places the non-blind signature layer on the corresponding sEUF-CMA footing. The public verifier bound
\(\BoundSig\) used by ARC-SPRUCE should therefore not be identified with \(\beta_{\mathrm{DKLW}}\) without a separate
norm-comparison lemma. In this section, \(\BoundSig\) is recorded only as the public acceptance bound of the concrete
verifier.

\subsection{SmallWood parameters}
SmallWood introduces parameters controlling packing, masking, and soundness:
\begin{itemize}[leftmargin=1.25em]
  \item the packing size \(s=|\Omega|\),
  \item the tail length \(\ell\) for HVZK masking and degree enforcement,
  \item the spot-check count \(\ell'\) and the number of masks \(\rho\) for PIOP soundness,
  \item the degree bounds for witness rows and constraint-derived masks, and
  \item the small-field knob \(\theta\) when using an extension-field variant.
\end{itemize}
The pre-sign proof is dominated by carrier membership, decode bridges, the inverse-witness check, and the
aligned public-target identity \(T=B_0+U\vecc+Z\). The post-sign proof adds the signature shortness gadget, the full
replay-image showing relation in eliminated BB-tran form, and the PRF checkpoint constraints.

\paragraph{Full replay-image geometry.}
Let \(n_{\mathrm{cols}}^{\mathrm{rep}}\) denote the row width used for the replay blocks, and write
\[
B_{\mathrm{rep}} := \left\lceil \frac{N}{n_{\mathrm{cols}}^{\mathrm{rep}}} \right\rceil .
\]
In the full replay-image showing model, each replayed ring object is represented by \(B_{\mathrm{rep}}\) replay blocks.
Accordingly, the post-sign witness carries
\[
4B_{\mathrm{rep}}
\]
replay rows for the full images of \(Au\), \(M+K\), \(R_0\), and \(R_1\), in addition to the cert-side signature
surface, source rows, and PRF companion rows.

\paragraph{Replay width as a proof-size knob.}
Increasing \(n_{\mathrm{cols}}^{\mathrm{rep}}\) decreases the replay block count \(B_{\mathrm{rep}}\) and hence the number
of replay-image rows and associated bridge families. This tends to reduce transcript size. The tradeoff is that larger
row width increases the degree of the committed row polynomials and therefore affects the LVCS/PCS soundness and opening
costs. The sizing discussion should therefore treat \(n_{\mathrm{cols}}^{\mathrm{rep}}\) as the main knob controlling the
transcript-size tradeoff of the full replay-image showing proof.

\paragraph{Replay-transform obligations for the showing theorem.}
The coefficient-domain consequence theorem of Section~\ref{subsec:statement-strength} relies on the exact full replay
transform \(\mathcal F_{\mathrm{rep}}\) and on the exact bridge equalities for the replay-image rows. In the full
replay-image model, injectivity is supplied by the transform itself rather than by an additional zero-detection
assumption on a reduced replay map; the replay residual transports the eliminated BB-tran relation directly.

\subsection{PRF parameters}
The PRF must be secure at the target security level, efficiently representable in constraints, and compatible with the
base field \(\Fq\). We use a Poseidon2-like permutation with feed-forward and truncation as a concrete instantiation.
Its parameters include the state width \(t\), round counts \((R_F,R_P)\), MDS matrices, round constants, S-box exponent
\(d\), and truncation length \(\ell_{\prftag}\). We require
\(\ell_{\prftag}\ge \lceil \lambda/\log_2(q)\rceil\) so the tag contains \(\lambda\) bits of entropy. The algorithm is
given in Section~\ref{subsec:poseidon-prf}, while the checkpointed constraint shape used in the post-sign proof is
stated in Section~\ref{subsec:prf-checkpointed}. Appendix~\ref{app:prf} records extended notes.

\subsection{Sizing discussion}
Concrete proof sizes depend jointly on the sampler parameters, the SmallWood masking parameters, the replay-block width,
and the PRF checkpoint schedule. In the full replay-image showing model, enlarging the replay-block width decreases the
number of committed replay rows, while the degree and opening costs of each row increase. Appendix~\ref{app:ext-params}
records representative parameter tuples and measurement data; the present section records only the qualitative coupling
between those choices.
